\documentclass[11pt]{article}
\input{../../shared/project_style.tex}
\rhead{Basics 19 Textbook Draft}

\begin{document}
\DocTitle{Basics 19: Plasma Distributions and Temperature}{III - Elementary plasma physics}{Velocity-space probability, Maxwellians, moments, anisotropy, and energetic tails}

\begin{overviewbox}
\textbf{Textbook draft, version 1.0.} A kinetic plasma is described by a distribution over position and velocity rather than one trajectory or one bulk velocity. This chapter builds the distribution-function language, derives Maxwellian moments, and clarifies the statistical meaning of plasma temperature. This chapter is designed for a reader with no prior plasma-physics training. It distinguishes exact identities from physical approximations and connects the central equations to later simulation modules.
\end{overviewbox}

\ProjectTOC

\section{Learning objectives}
After completing this note, the reader should be able to:
\begin{itemize}[leftmargin=1.6em]
\item Interpret the distribution function as phase-space density.
\item Normalize one- and three-dimensional Maxwellian distributions.
\item Derive density, flow velocity, pressure, and temperature moments.
\item Distinguish common thermal-speed conventions.
\item Represent anisotropic and nonthermal populations.
\item Connect particle-marker samples to continuum moments.
\end{itemize}

\section{Prerequisites and reading strategy}
\begin{itemize}[leftmargin=1.6em]
\item Basics 1, 2, 6, 10, and 17.
\end{itemize}
On a first reading, emphasize physical meaning and dimensional checks. On a second reading, reproduce the derivations. Attempt the computational laboratory only after the limiting cases are understood.

\section{Distribution function and normalization}

For species $s$, let $f_s(\mathbf x,\mathbf v,t)$ satisfy
\begin{equation}
dN_s=f_s\,d^3x\,d^3v.
\end{equation}
Its velocity integral is number density:
\begin{equation}
n_s(\mathbf x,t)=\int f_s\,d^3v.
\end{equation}
Some communities normalize $f$ to unity and write $nf$ for phase-space density. The convention must be explicit because it changes every moment formula.


\section{One-dimensional Maxwellian}

Thermal equilibrium in one velocity component has
\begin{equation}
f(v_x)=n\sqrt{\frac{m}{2\pi k_BT}}
\exp\left[-\frac{m(v_x-u_x)^2}{2k_BT}\right].
\end{equation}
The Gaussian integral ensures $\int f\,dv_x=n$. Its mean is $u_x$, and its variance is
\begin{equation}
\left\langle(v_x-u_x)^2\right\rangle=\frac{k_BT}{m}.
\end{equation}
Temperature therefore measures random kinetic-energy spread, not bulk flow speed.


\section{Three-dimensional Maxwellian and kinetic energy}

For an isotropic equilibrium,
\begin{equation}
f_M(\mathbf v)=n\left(\frac{m}{2\pi k_BT}\right)^{3/2}
\exp\left[-\frac{m|\mathbf v-\mathbf u|^2}{2k_BT}\right].
\end{equation}
Each quadratic degree of freedom contributes $k_BT/2$:
\begin{equation}
\frac{m}{2n}\int|\mathbf v-\mathbf u|^2f_M\,d^3v=\frac{3}{2}k_BT.
\end{equation}
Bulk kinetic energy $mn u^2/2$ is separate from internal thermal energy.


\section{Moments and the pressure tensor}

The bulk velocity is
\begin{equation}
\mathbf u_s=\frac{1}{n_s}\int\mathbf v f_s\,d^3v.
\end{equation}
Define peculiar velocity $\mathbf c=\mathbf v-\mathbf u_s$. The pressure tensor is
\begin{equation}
\mathsf P_s=m_s\int\mathbf c\mathbf c\,f_s\,d^3v.
\end{equation}
For an isotropic distribution, $\mathsf P_s=p_s\mathsf I$ and
\begin{equation}
p_s=n_sk_BT_s=\frac{1}{3}\operatorname{tr}\mathsf P_s.
\end{equation}
The heat flux is a third-order moment,
\begin{equation}
\mathbf q_s=\frac{m_s}{2}\int c^2\mathbf c f_s\,d^3v.
\end{equation}


\section{Thermal-speed conventions}

Several definitions are common:
\begin{equation}
v_T=\sqrt{\frac{k_BT}{m}},\qquad
v_{\rm mp}=\sqrt{\frac{2k_BT}{m}},\qquad
v_{\rm rms}=\sqrt{\frac{3k_BT}{m}}.
\end{equation}
They are respectively a one-component standard deviation, most-probable speed, and root-mean-square speed. Calling all three ``thermal speed'' without a definition creates factor-of-$\sqrt{2}$ or $\sqrt{3}$ errors.


\section{Anisotropic distributions}

Magnetized collisionless plasmas often have different parallel and perpendicular temperatures. A bi-Maxwellian can be written
\begin{equation}
f_b=
n\left(\frac{m}{2\pi k_BT_\perp}\right)
\left(\frac{m}{2\pi k_BT_\parallel}\right)^{1/2}
\exp\left[-\frac{mv_\perp^2}{2k_BT_\perp}
-\frac{m(v_\parallel-u_\parallel)^2}{2k_BT_\parallel}\right].
\end{equation}
Its pressure tensor has $p_\parallel=nk_BT_\parallel$ and $p_\perp=nk_BT_\perp$. Such anisotropy can drive kinetic instabilities.


\section{Nonthermal and energetic populations}

Fusion alpha particles and neutral-beam ions are not generally Maxwellian. A slowing-down distribution has a characteristic form
\begin{equation}
f_{\rm sd}(v)\propto\frac{1}{v^3+v_c^3}
\end{equation}
over a finite birth-energy range. Beams may be strongly anisotropic. A small energetic population can dominate wave drive even when its density is much lower than the thermal-plasma density because resonance depends on phase-space gradients and particle energy.


\section{Particle markers and sampling error}

A Monte Carlo representation uses weighted markers:
\begin{equation}
f(\mathbf z)\approx\sum_{p=1}^{N_p}w_pK_h(\mathbf z-\mathbf z_p).
\end{equation}
Moments become weighted sums. Statistical error typically decreases as $N_p^{-1/2}$ for independent samples, while interpolation kernels introduce bias and smoothing. Verification should compare sampled moments with analytic distribution moments.


\section{Computational laboratory}
Sample one- and three-dimensional Maxwellians, reconstruct histograms, and estimate density, flow, pressure, and temperature. Compare thermal-speed definitions. Repeat for a bi-Maxwellian and a simple slowing-down distribution, then quantify Monte Carlo convergence.

\section{Summary}
\begin{itemize}[leftmargin=1.6em]
\item A distribution function describes particle density in phase space.
\item Temperature is a moment of random velocity spread, not bulk motion.
\item Pressure is generally a tensor and becomes scalar only under isotropy.
\item Thermal-speed conventions must be stated explicitly.
\item Energetic and anisotropic populations require non-Maxwellian models and careful sampling.
\end{itemize}

\SymbolsAcronymsSection
\begin{symtable}
$f_s$ & distribution function & Species phase-space density. \\
$\mathbf u_s$ & bulk velocity & First velocity moment divided by density. \\
$\mathsf P_s$ & pressure tensor & Second central velocity moment. \\
$\mathbf q_s$ & heat flux & Third central moment. \\
$\mathbf c$ & peculiar velocity & $\mathbf v-\mathbf u_s$. \\
$T_\parallel,T_\perp$ & anisotropic temperatures & Parallel and perpendicular random-energy measures. \\
\end{symtable}

\section{Exercises}
\begin{enumerate}[leftmargin=1.8em]
\item Normalize the one-dimensional Maxwellian.
\item Derive $p=nk_BT$ from the isotropic pressure tensor.
\item Separate bulk and thermal kinetic energy for a shifted Maxwellian.
\item Calculate $T_\parallel$ and $T_\perp$ from sampled velocities.
\item Explain why a low-density energetic tail can strongly affect an Alfvén eigenmode.
\end{enumerate}

\section{References}
\begin{thebibliography}{99}
\bibitem{ref1} G. B. Arfken, H. J. Weber, and F. E. Harris, \emph{Mathematical Methods for Physicists}, 7th ed., Academic Press (2013).
\bibitem{ref2} G. Strang, \emph{Linear Algebra and Its Applications}, 4th ed., Brooks/Cole (2006).
\bibitem{ref3} W. A. Strauss, \emph{Partial Differential Equations: An Introduction}, 2nd ed., Wiley (2007).
\bibitem{ref4} F. F. Chen, \emph{Introduction to Plasma Physics and Controlled Fusion}, 3rd ed., Springer (2016).
\bibitem{ref5} R. J. Goldston and P. H. Rutherford, \emph{Introduction to Plasma Physics}, CRC Press (1995).
\bibitem{ref6} P. M. Bellan, \emph{Fundamentals of Plasma Physics}, Cambridge University Press (2006).
\bibitem{ref7} T. H. Stix, \emph{Waves in Plasmas}, American Institute of Physics (1992).
\bibitem{ref8} D. A. Gurnett and A. Bhattacharjee, \emph{Introduction to Plasma Physics}, 2nd ed., Cambridge University Press (2017).
\end{thebibliography}

\end{document}
